%% Manuscript draft. Compile with pdflatex + bibtex.
%% Numbers kept in sync with results/tables/*.json and *.csv (see manuscript/README.md).
\documentclass[10pt,twocolumn]{article}
\usepackage[margin=0.75in]{geometry}
\usepackage{amsmath,amssymb,amsthm,mathtools}
\usepackage{graphicx,booktabs,multirow}
\usepackage[colorlinks=true,linkcolor=blue,citecolor=blue]{hyperref}
\usepackage{natbib}
\bibliographystyle{plainnat}

\newtheorem{theorem}{Theorem}
\newtheorem{lemma}{Lemma}
\newtheorem{corollary}{Corollary}
\newtheorem{proposition}{Proposition}
\newtheorem{assumption}{Assumption}
\newtheorem{remark}{Remark}

\newcommand{\E}{\mathbb{E}}
\newcommand{\R}{\mathbb{R}}
\newcommand{\phidi}{\phi_{\mathrm{DI}}}
\newcommand{\phidm}{\phi_{\mathrm{DM}}}
\newcommand{\dhat}{\hat\Delta}
\newcommand{\tbar}{\bar\theta}
\newcommand{\Dp}{D'}
\newcommand{\Dt}{\tilde D}
\newcommand{\todo}[1]{\textbf{[TODO: #1]}}

\title{Separating Training, Evaluation and Constraint-Set Noise\\in Fairness Uncertainty Quantification}
\author{\todo{author list}}
\date{\today}

\begin{document}
\maketitle

\begin{abstract}
Fairness-aware classifiers trained by stochastic gradient methods are random objects, so any reported
demographic-parity or equalised-odds gap is random as well.  Recent work constructs confidence
intervals for such gaps with an online multiplier bootstrap, but those intervals quantify only the
randomness of the optimisation trajectory.  With the held-out sets used in practice, a second and
often larger component---the sampling noise of the evaluation set---is ignored, and a third arises
from the random constraint set that defines the fairness-aware estimator.  We decompose the error of
the reported gap into these asymptotically independent parts, derive the limiting distribution and
the exact coverage loss of the existing interval as a function of the ratio between training
iterations and evaluation size, and prove that it is valid only when the evaluation set is large
relative to the training run.  We then propose a joint bootstrap that resamples all three sources,
prove its asymptotic validity and give a sample-size allocation rule.  In a reproduction on synthetic
data, Adult, ACSIncome and COMPAS, the existing interval covers the population fairness gap with
probability between $0.00$ and $0.62$, whereas the joint interval attains nominal coverage (on Adult,
$0.98$ versus $0.30$).  A variance budget shows that training noise accounts for only
$0.8\%$--$8.3\%$ of the total uncertainty, with evaluation and constraint-set sampling dominating.

\par\medskip\noindent\textbf{Keywords:} fairness auditing; uncertainty quantification; constrained stochastic optimization; online bootstrap; distribution-free inference.
\end{abstract}

\section{Introduction}
\label{sec:intro}

Machine-learning models used in high-stakes decisions are increasingly required to report fairness
diagnostics such as the demographic-parity (DP) gap or the equalised-odds (EO) gap between protected
groups \citep{hardt2016,corbett2017}.  In practice these diagnostics are computed from a finite
held-out sample for a classifier whose parameters were themselves produced by a stochastic
optimisation algorithm.  Both steps introduce randomness, and a point estimate of a fairness gap
without a measure of uncertainty is therefore of limited use for auditing.

\citet{roy2026uq} take up this problem for fairness-aware linear classifiers trained online by
stochastic dual averaging (for DP) or penalised SGD (for EO-type criteria).  Viewing training as
constrained stochastic optimisation, they establish asymptotic normality of the averaged iterate, and
because the asymptotic covariance depends on the unknown optimum they propose an online multiplier
bootstrap operating on a single data stream.  The resulting interval is appealing: it is
computationally cheap, parallelisable, and requires no repeated training.

This paper starts from a simple observation about what such an interval actually covers.  Write
$\theta^*$ for the population optimum of the fairness-aware objective, $\tbar_T$ for the averaged
online iterate after $T$ steps, and $\dhat_m$ for the empirical fairness gap computed on an
independent held-out set $\Dt$ of size $m$.  The reported number is $\dhat_m(\tbar_T)$, while the
quantity a deployment decision depends on is the population gap $\Delta(\theta^*)$.  The error
splits into
\[
\dhat_m(\tbar_T)-\Delta(\theta^*)
=\underbrace{\Delta(\tbar_T)-\Delta(\theta^*)}_{\text{training}}
+\underbrace{\dhat_m(\tbar_T)-\Delta(\tbar_T)}_{\text{evaluation}} ,
\]
and the two terms are independent because $\Dt$ is independent of the training stream.
The online bootstrap approximates the law of the first term only.  If the evaluation error is small
relative to the training noise the distinction is immaterial; in the operating regime of interest it
is not.  In our reproduction, with $m=1000$ held-out observations, the evaluation noise for Adult is
$3.9$ times the training noise, and the existing interval covers the population gap in $30\%$ of
repetitions instead of the nominal $95\%$.

\paragraph{Contributions.}
\begin{enumerate}
  \item \emph{Decomposition and a diagnostic.}  We show that the error of the reported gap decomposes
        into three asymptotically independent components---training, evaluation, and constraint-set
        sampling---and derive the limiting distribution (Theorem~\ref{thm:limit}) together with the
        exact coverage of the existing interval (Corollary~\ref{cor:coverage}).  The interval is
        asymptotically valid only when $m \gg T\,v^2/\sigma^2_\Delta$, a condition violated by orders
        of magnitude in published experimental setups.
  \item \emph{A joint bootstrap.}  We propose to combine the multiplier bootstrap over training
        trajectories with a resampling of the evaluation set, and in addition an outer bootstrap over
        the constraint set that defines $\theta^*$; the construction is asymptotically valid
        (Theorem~\ref{thm:joint}) and costs one extra pass over the evaluation set, plus a small
        number of direct constrained solves for the constraint-set component.
  \item \emph{A variance budget and design guidance.}  On Adult, ACSIncome and COMPAS the training
        component accounts for only $0.8$--$8.3\%$ of the total variance of a reported gap; the
        dominant components are evaluation sampling ($37$--$78\%$) and constraint-set sampling
        ($18$--$55\%$).  We translate this into a rule for allocating audit effort between training
        iterations, labelled evaluation data and the choice of constraint set.
  \item \emph{A reproducible study.}  We re-implement the baseline from the published equations
        (no official code was released), reproduce its qualitative claims on five data settings, and
        document two further reproducibility findings: the published constant step size does not
        reach the reported convergence regime, and the synthetic data generator is close to
        degenerate.  All code, data splits and raw results are released.
\end{enumerate}

The remainder of the paper is organised as follows.  Section~\ref{sec:related} reviews related work,
Section~\ref{sec:setup} fixes notation, Section~\ref{sec:decomp} states the decomposition,
Section~\ref{sec:method} describes the joint bootstrap, Section~\ref{sec:exp} reports experiments and
Section~\ref{sec:discussion} discusses limitations and practical implications.

\section{Related work}
\label{sec:related}

\paragraph{Online inference for stochastic approximation.}
Asymptotic normality of averaged stochastic approximation goes back to
\citet{polyak1992}; constrained variants are analysed by \citet{duchi2016}.  Online bootstrap and
batch-means procedures for unconstrained SGD are developed by \citet{fang2018,chen2016,zhu2021} and
extended to policy evaluation \citep{ramprasad2023}.  These methods target the model parameter; for
fairness we need the law of a non-smooth functional of it, which is why \citet{roy2026uq} combine a
functional delta-method argument with the multiplier bootstrap.  Our work does not modify the
training-side theory; it adds the two missing sources of randomness on top of it.

\paragraph{Statistical inference for fairness.}
A parallel literature tests whether a \emph{fixed} classifier is fair: permutation tests
\citep{diciccio2020}, optimal-transport projections \citep{taskesen2021}, adversarial individual-fairness
tests \citep{maity2021}, and bootstrap audits over rich group collections \citep{cherian2023}.  These
papers treat the classifier as given and quantify evaluation uncertainty only, which is the
complementary half of our decomposition.  Our method targets the composition of training and
evaluation uncertainty and makes the two components separately reportable.

\paragraph{Uncertainty in fairness more broadly.}
\citet{rosenblatt2024} benchmark uncertainty estimates used \emph{inside} fair predictors, a different
object from uncertainty about the fairness of a trained predictor.  \citet{moreau2024} argue
conceptually for uncertainty-aware fairness reporting.  Work on the variance of subgroup-level
fairness estimates has highlighted that minority groups carry disproportionate estimation variance;
our evaluation component formalises this in the DI/DM setting and is combined with the training
component rather than studied in isolation.

\paragraph{Reproducibility and its use here.} Neither the published paper nor its conference version \citep{roy2023uq} ships official code, so we
re-implemented both training procedures and the multiplier bootstrap from the published equations and
report a full reproduction study.  The qualitative claims reproduce: the mean interval score
decreases with the number of iterations, the empirical distribution of the estimated gap converges to
the benchmark normal law in Wasserstein-2 distance, and the interval covers the held-out estimand.
Two discrepancies are documented: the published constant step size ($\eta_k=0.01k^{-0.501}$) does not
reach the convergence regime displayed in the paper within $10^6$ iterations, and the synthetic
generator is close to degenerate because the protected attribute is nearly a copy of the label.  We
therefore calibrate the step size (and, on real data, the regularisation strength) and report both
the paper-literal and the calibrated settings.

\section{Problem setup}
\label{sec:setup}

\paragraph{Data and classifier.} Let $(x,y,z)\sim P$ with features $x\in\R^p$, label
$y\in\{-1,1\}$ and protected attribute $z\in\{0,1\}$.  We consider linear classifiers
$h(x;\theta)=\operatorname{sign}(\theta^\top\varphi(x))$ with $\varphi(x)=(x^\top,1)^\top
\in\R^d$, $d=p+1$.

\paragraph{Disparate impact.} With tolerance $\epsilon>0$ and a held-out constraint set
$\Dp$ of size $n_c$, define
$\tilde x=n_c^{-1}\sum_{i\in\Dp}(z_i-\bar z)\varphi(x_i)$ and
$\Theta(\Dp)=\{\theta:|\tilde x^\top\theta|\le\epsilon\}$.  The target parameter is
\begin{equation}\label{eq:theta}
\theta^*=\arg\min_{\theta\in\Theta(\Dp)}\ \ell(\theta),\qquad
\ell(\theta)=\E[\log(1+e^{-y\theta^\top\varphi(x)})]+\frac{\kappa}{2}\|\theta\|_2^2 ,
\end{equation}
and the population gap is $\Delta(\theta)=P(h=1\mid z=0)-P(h=1\mid z=1)$ with
$\phi(\theta)=|\Delta(\theta)|$.

\paragraph{Disparate mistreatment.} With penalty $R^2>0$ and $\Dp_-=\{i\in\Dp:y_i=-1\}$,
$\gamma(\theta;\Dp_-)=\big(n_-^{-1}\sum_{i\in\Dp_-}(z_i-\bar z)\min(0,y_i\varphi(x_i)^\top
\theta)\big)^2$ and $\theta^*=\arg\min_\theta \ell(\theta)+R^2\gamma(\theta;\Dp_-)$; the population
gap conditions on $y=-1$.

\paragraph{Online training.} The averaged iterate $\tbar_T=T^{-1}\sum_{k=1}^T\theta_k$ is produced by
stochastic dual averaging with projection onto $\Theta(\Dp)$ for DP, or by penalised SGD for EO.  An
online multiplier bootstrap with $B$ perturbed trajectories, each driven by the same single data
stream with i.i.d.\ multipliers $V\sim\mathrm{Unif}[1-\sqrt3,1+\sqrt3]$, yields
$\{\hat\theta^b_T\}_{b=1}^B$ and hence $\{\dhat_m(\bar\theta^b_T)\}$; see \citet{roy2026uq} for the
algorithm and its training-side guarantees.

\paragraph{Two estimands.} Notation is collected in Table~\ref{tab:notation}.  We distinguish the \emph{held-out estimand} $\dhat_m(\theta^*)$, i.e.\ the
fairness of the optimal classifier evaluated on the very sample used by the auditor, from the
\emph{population estimand} $\Delta(\theta^*)$, i.e.\ the fairness of the deployed classifier on the
data distribution.  The distinction is vacuous only when $m\to\infty$; a confidence interval must
name its target.

\begin{table}[t]
\centering
\small
\caption{Notation.}
\label{tab:notation}
\resizebox{\linewidth}{!}{\begin{tabular}{ll}
\toprule
symbol & meaning \\
\midrule
$\theta$ & linear-classifier parameter; $h(x;\theta)=\operatorname{sign}(\theta^\top\varphi(x))$ \\
$\theta^*$ & population constrained optimum \eqref{eq:theta} \\
$\tbar_T$ & averaged online iterate after $T$ steps \\
$\Dp$, $n_c$ & constraint set (audit design) and its size \\
$\Dt$, $m$ & held-out evaluation set and its size \\
$\Delta(\theta)$, $\dhat_m(\theta)$ & population and empirical fairness gap \\
$\sigma^2_\Delta$, $v^2$ & training and evaluation variance constants \\
$\lambda=T/m$ & training-to-evaluation effort ratio \\
$\sigma^2_{\rm tot}$ & $\sigma^2_\Delta+\lambda v^2$, limiting variance of the reported gap \\
$B$ & number of multiplier-bootstrap trajectories \\
$\eta_k$, $\kappa$, $R$ & step size, $L_2$ regularisation, mistreatment penalty \\
\bottomrule
\end{tabular}}
\end{table}

\section{Three-source decomposition}
\label{sec:decomp}

\begin{assumption}[inherited from \citealp{roy2026uq}]
\label{ass:base}
The loss is strongly convex and smooth on a neighbourhood of $\theta^*$; the data stream is i.i.d.;
the decision boundary is linear in $\theta$; the conditional feature distributions admit continuous
densities at the boundary.  The evaluation set $\Dt$ is drawn i.i.d.\ from $P$ independently of the
training stream and of $\Dp$.
\end{assumption}

\begin{lemma}[training noise; \citealp{roy2026uq}]
\label{lem:train}
Under Assumption~\ref{ass:base},
$\sqrt T(\Delta(\tbar_T)-\Delta(\theta^*))\Rightarrow\mathcal N(0,\sigma^2_\Delta)$ with
$\sigma^2_\Delta=g^\top\Sigma g$, $g=\nabla\Delta(\theta^*)$, and $\Sigma$ the asymptotic covariance
of the averaged iterate.
\end{lemma}

\begin{lemma}[evaluation noise]
\label{lem:eval}
Under Assumption~\ref{ass:base}, uniformly over $\theta$ in a neighbourhood of $\theta^*$,
\[
\sqrt m\big(\dhat_m(\theta)-\Delta(\theta)\big)\Rightarrow\mathcal N\!\big(0,\,v^2(\theta)\big),
\qquad
v^2(\theta)=\frac{p_0(1-p_0)}{P(z=0)}+\frac{p_1(1-p_1)}{P(z=1)},
\]
with $p_j=P(h=1\mid z=j)$, and the analogous expression with $y=-1$ for DM.  Moreover the two
convergences are independent.
\end{lemma}

\begin{theorem}[limit distribution]
\label{thm:limit}
Suppose $\Delta(\theta^*)\neq0$ and $T/m\to\lambda\in[0,\infty)$.  Then under
Assumption~\ref{ass:base},
\[
\sqrt T\big(\dhat_m(\tbar_T)-\Delta(\theta^*)\big)\Rightarrow
\mathcal N\big(0,\ \sigma^2_\Delta+\lambda v^2\big).
\]
\end{theorem}

\begin{corollary}[coverage of the existing interval]
\label{cor:coverage}
The fixed-time interval of \citet{roy2026uq}, of asymptotic half-width
$z_{\alpha/2}\sigma_\Delta/\sqrt T$, has asymptotic coverage
\[
2\Phi\!\Big(z_{\alpha/2}\,\frac{\sigma_\Delta}{\sqrt{\sigma^2_\Delta+\lambda v^2}}\Big)-1,
\]
which equals $1-\alpha$ if and only if $\lambda v^2/\sigma_\Delta^2\to0$, i.e.\
$m\gg T v^2/\sigma_\Delta^2$.
\end{corollary}

\begin{remark}[two coverage protocols]
\label{rem:protocol}
If the evaluation set $\Dt$ is held fixed across repetitions, as in the original simulation design,
the evaluation error becomes a constant bias $b=\dhat_m(\theta^*)-\Delta(\theta^*)$ and coverage is
$\Phi(z_{\alpha/2}-|b|\sqrt T/\sigma_\Delta)-\Phi(-z_{\alpha/2}-|b|\sqrt T/\sigma_\Delta)$.
Coverage statements about population fairness must therefore be evaluated with a fresh evaluation
set in each repetition; we report both protocols.
\end{remark}

\section{Method: joint bootstrap}
\label{sec:method}

For $b=1,\dots,B$ multiplier trajectories, $r=1,\dots,R$ evaluation resamples, and $s=1,\dots,S$
constraint-set resamples, define the bootstrap deviations
\[
D^*_{b,r,s}
= \underbrace{\dhat_m(\bar\theta^b_T)-\dhat_m(\tbar_T)}_{\text{training}}
+ \underbrace{\dhat_{m,r}(\tbar_T)-\dhat_m(\tbar_T)}_{\text{evaluation}}
+ \underbrace{\dhat_{m}(\theta^*_s)-\dhat_m(\theta^*)}_{\text{constraint set}} ,
\]
where the third term uses the constrained optimum $\theta^*_s$ computed for a resampled constraint
set $\Dp_s$; it is obtained by solving the constrained problem directly, with no retraining.
The proposed interval is
$\big[\dhat_m(\tbar_T)+\hat q_{\alpha/2},\ \dhat_m(\tbar_T)+\hat q_{1-\alpha/2}\big]$
for the quantiles of $D^*$, mapped to $|\Delta|$ by folding: if the resulting interval for the signed
gap excludes zero we report the absolute values of its endpoints, otherwise we report
$[0,\max(|l|,|u|)]$.  When the constraint set is regarded as fixed by design, the third term is
omitted and the interval targets $\Delta(\theta^*)$ for the realised $\Dp$.

\begin{theorem}[asymptotic validity]
\label{thm:joint}
Under Assumption~\ref{ass:base}, if $S\to\infty$ and $\tbar_T\to\theta^*$, then conditionally on the
data the law of $D^*$ converges in probability to the law of
$\dhat_m(\tbar_T)-\Delta(\theta^*)$; consequently the percentile interval has asymptotic coverage at
least $1-\alpha$ for $\phi(\theta^*)$.
\end{theorem}

\paragraph{Cost.} The evaluation resampling is a vectorised operation on the $m$ held-out
predictions and costs $O(Rm)$.  The training deviation reuses the $B$ trajectories already produced
by the online bootstrap.  The constraint-set bootstrap requires $S$ direct constrained solves, each
costing a few iterations of a quasi-Newton method, and no additional passes over the data stream.

\paragraph{Sample-size allocation.}
\label{sec:ss}
With per-unit costs $c_T$ and $c_m$ and target half-width $w$, minimising $c_T T+c_m m$ subject to
$z_{\alpha/2}^2(\sigma^2_\Delta/T+v^2/m)\le w^2$ gives
\[
\frac{m}{T}=\sqrt{\frac{v^2c_T}{\sigma^2_\Delta c_m}},\qquad
T=\frac{z_{\alpha/2}^2}{w^2}\Big(\sigma^2_\Delta+\sqrt{\sigma^2_\Delta v^2 c_T/c_m}\Big).
\]
Because $\sigma^2_\Delta$ and $v^2$ are of the same order in the regimes we study, the binding
constraint is the \emph{absolute} size of $m$: increasing $T$ alone cannot push the half-width below
$z_{\alpha/2}v/\sqrt m$ (Figure~\ref{fig:samplesize}).

\section{Experiments}
\label{sec:exp}

\paragraph{Setup and protocols.} Synthetic data follow the generators of \citet{zafar2019} used by
\citet{roy2026uq}; real data are Adult and ACSIncome \citep{ding2021} and COMPAS, with the sensitive
attributes and tolerances of the original Table~1.  Three protocols appear below.
\emph{(R) Reproduction}: one fixed constraint set and one fixed held-out set, repetitions resample
only the training stream, which is the protocol of the original simulation study.
\emph{(A) Protocol A}: the constraint set is fixed but every repetition draws a fresh evaluation set
from a disjoint reservoir, so that coverage is evaluated over the evaluation sampling law as well.
\emph{(B) Budget}: a single run with the reported audit budget ($T=3\times10^4$, $m=1000$ for
Adult/ACSIncome and $m=400$ for COMPAS), used for the sample-size analysis.  Unless stated otherwise
$B=100$ multiplier trajectories and $\alpha=0.05$; step size and regularisation are calibrated as
described above ($\eta_0=0.1$ for DI, $0.03$ for the synthetic DM generator, $\kappa=0.1$).

\paragraph{Reproduction of the baseline.}
Table~\ref{tab:repro} summarises the reproduced behaviour in the original fixed-held-out protocol.
The qualitative claims of the baseline hold: the mean interval score decreases with $T$, the
empirical distribution converges to the benchmark normal law, and the interval covers the held-out
estimand.  For the population estimand, however, coverage ranges from $0.00$ to $0.62$ except in the
synthetic DM setting, where the interval is wide enough by accident.  Figures~\ref{fig:reprodi}
and~\ref{fig:reprocompas} show the per-iteration trajectories behind these numbers: convergence to
the asymptotic law is fast on the synthetic experiment, the interval covers the held-out estimand
throughout, and the population estimand is missed as soon as the evaluation error is comparable to
the interval width. The $T^{-1/2}$ rate behind Theorem~\ref{thm:limit} is itself visible in a separate grid of training horizons: the distance of $\tbar_T$ to the constrained optimum decays with log--log slope $0.47$ and the signed training bias of the gap with slope $0.46$ (see the empirical convergence rates in the supplementary material).

\begin{figure}[t]
\centering
\includegraphics[width=\linewidth]{../results/figures/repro_synth_di.pdf}
\caption{Reproduction of the synthetic DI experiment: (a) Wasserstein distance between the empirical
law of $\dhat_m(\tbar_T)$ and its asymptotic limit; (b) mean interval score; (c) coverage of the
held-out estimand $\dhat_m(\theta^*)$ and of the population estimand $\Delta(\theta^*)$;
(d) interval width against the sampling spread $4v/\sqrt m$.  The interval covers the held-out
estimand but not the population estimand, the failure mechanism analysed in
Section~\ref{sec:decomp}.}
\label{fig:reprodi}
\end{figure}

\begin{figure}[t]
\centering
\includegraphics[width=\linewidth]{../results/figures/repro_compas.pdf}
\caption{Reproduction of the COMPAS DM experiment (same four panels).  The penalty term is frozen at
its training value inside the feasible region, so the averaged iterate stops short of the
constrained boundary; the held-out estimand is covered because the evaluation term dominates, while
the population estimand is covered in only $0.62$ of repetitions.}
\label{fig:reprocompas}
\end{figure}

\begin{table}[t]
\centering
\small
\setlength{\tabcolsep}{4pt}
\caption{Reproduction of the baseline in the fixed-held-out protocol.  ``coverage held-out'' is the
fraction of repetitions whose interval contains $\dhat_m(\theta^*)$; ``coverage population'' uses
$\Delta(\theta^*)$ evaluated on the full data.}
\label{tab:repro}
\resizebox{\linewidth}{!}{\begin{tabular}{lcccccc}
\toprule
setting & $T$ & reps & $\phi(\theta^*)$ & cov.\ held-out & cov.\ population & width \\
\midrule
Synthetic (DI) & 2e4 & 100 & 0.0835 & 0.89 & 0.43 & 0.077 \\
Synthetic (DM) & 2e4 & 100 & 0.1401 & 0.86 & 1.00 & 0.208 \\
Adult (DI)     & 5e4 & 100 & 0.0564 & 1.00 & 0.00 & 0.022 \\
ACSIncome (DI) & 5e4 & 150 & 0.0442 & 1.00 & 0.00 & 0.013 \\
COMPAS (DM)    & 5e4 & 150 & 0.0133 & 1.00 & 0.62 & 0.115 \\
\bottomrule
\end{tabular}}
\end{table}

\paragraph{Protocol-A coverage.}
Table~\ref{tab:protocolA} varies the held-out size $m$ with a fresh evaluation set per repetition.
The baseline coverage increases with $m$ exactly as predicted by Corollary~\ref{cor:coverage}, while
the joint bootstrap stays at or above nominal coverage; its width converges to that of the baseline
as $m$ grows, quantifying the price of valid inference.  The same pattern holds on the three real
datasets, where the baseline interval attains $0.25$--$0.62$ coverage and the joint interval
$0.98$--$1.00$.

\begin{table}[t]
\centering
\small
\setlength{\tabcolsep}{4pt}
\caption{Protocol-A coverage of the population gap (nominal $0.95$) with a fresh evaluation set in
every repetition.  Synthetic DI: $T=2\times10^4$, $B=100$, 100 repetitions,
$\phi(\theta^*)=0.0212$.  Real data: $T=3\times10^4$, $B=100$, 60 repetitions.}
\label{tab:protocolA}
\resizebox{\linewidth}{!}{\begin{tabular}{lcccccc}
\toprule
setting & $m$ & cov.\ base & cov.\ joint & $w$ base & $w$ joint & ratio \\
\midrule
Synthetic (DI) & 100  & 0.56 & 0.99 & 0.107 & 0.283 & 2.64 \\
Synthetic (DI) & 1000 & 0.85 & 0.98 & 0.077 & 0.116 & 1.50 \\
Synthetic (DI) & 5000 & 0.99 & 1.00 & 0.074 & 0.085 & 1.15 \\
Adult (DI)     & 1000 & 0.30 & 0.98 & 0.025 & 0.102 & 4.15 \\
ACSIncome (DI) & 1000 & 0.25 & 1.00 & 0.015 & 0.082 & 5.72 \\
COMPAS (DM)    & 400  & 0.62 & 1.00 & 0.096 & 0.173 & 2.08 \\
\bottomrule
\end{tabular}}
\end{table}

\paragraph{Variance budget.}
Table~\ref{tab:budget} decomposes the variance of the reported gap.  Training noise, the only
component addressed by the baseline, contributes between $0.8\%$ and $8.3\%$ of the total.  The
constraint-set component is particularly severe for COMPAS, where resampling $\Dp$ changes the sign
of the reported gap.

\begin{table}[t]
\centering
\small
\setlength{\tabcolsep}{4pt}
\caption{Variance budget for the reported population gap: standard errors of the three components and
their shares of the total variance.}
\label{tab:budget}
\resizebox{\linewidth}{!}{\begin{tabular}{lcccccc}
\toprule
dataset & se(train) & se(eval) & se($\Dp$) & share(train) & share(eval) & share($\Dp$) \\
\midrule
Adult (DI)     & 0.0056 & 0.0259 & 0.0124 & 3.6\% & 78.3\% & 18.1\% \\
ACSIncome (DI) & 0.0034 & 0.0304 & 0.0245 & 0.8\% & 60.3\% & 39.0\% \\
COMPAS (DM)    & 0.0293 & 0.0621 & 0.0753 & 8.3\% & 37.1\% & 54.6\% \\
\bottomrule
\end{tabular}}
\end{table}

\paragraph{Sample-size planning and the minimum certifiable gap.}
Table~\ref{tab:budget} and the planning formula of Section~\ref{sec:ss} together answer a practical
question: at a given audit budget, which fairness gaps can be certified at all?  Using the component
estimates from the protocol-A runs, the best achievable $95\%$ half-width at $T=3\times10^4$ and
$m=1000$ is $0.053$ for Adult, $0.060$ for ACSIncome and $0.133$ for COMPAS ($m=400$).  The reported
population gaps are $0.059$, $0.041$ and $0.006$ respectively, so at these budgets the audit can
exclude a zero gap for Adult only; for ACSIncome and COMPAS the minimum certifiable gap exceeds the
gap itself, by a factor of $1.5$ and $21$, respectively.  Figure~\ref{fig:samplesize} shows the
feasible $(T,m)$ region for several target half-widths.  The practical implication is that increasing
$T$ alone is not enough: because $\sigma_\Delta^2$ and $v^2$ are of the same order, the half-width
floor $z_{\alpha/2}v/\sqrt m$ sets a hard limit that only more evaluation data can lower.

\begin{figure}[t]
\centering
\includegraphics[width=\linewidth]{../results/figures/sample_size.pdf}
\caption{Feasible $(T,m)$ region for target $95\%$ half-widths (curves), with the audit budget
($\star$) and the reported population gap annotated in each panel.  Points below the curve for a
target width cannot certify that width.}
\label{fig:samplesize}
\end{figure}

\paragraph{Optimisation bias.}
\label{par:bias}
A variance-based interval presumes that the averaged iterate has reached the neighbourhood of
$\theta^*$.  Table~\ref{tab:sensitivity} varies the regularisation strength on Adult under protocol
(R).  At $\kappa=10^{-4}$ the optimum has norm $10.9$ and the averaged iterate is $10.5$ away from
it: no interval based on the training variance can be valid.  The table also shows that coverage is
governed by the \emph{functional} bias of $\phi$, not by the parameter distance: at $\kappa=10^{-2}$
the parameter bias is $0.81$ yet coverage is nominal, while at $\kappa=0.1$ a parameter bias of
$0.12$ is enough to lose all coverage because the DI functional shifts by more than the half-width
(DP/DM gaps are discontinuous in $\theta$).  Audits should therefore report a convergence diagnostic
together with the interval, and the optimisation bias must be controlled by the design ($T$ and
$\kappa$), not by the bootstrap.

\begin{table}[t]
\centering
\caption{Optimisation-bias sensitivity on Adult (protocol R, 20 repetitions): the averaged iterate's
distance to the constrained optimum, the population coverage of the baseline interval, and its width.}
\label{tab:sensitivity}
\resizebox{\linewidth}{!}{\begin{tabular}{cccccc}
\toprule
$\kappa$ & $\|\theta^*\|$ & median $\|\tbar_T-\theta^*\|$ & coverage & CI width & $\phi(\theta^*)$ \\
\midrule
$10^{-4}$ & 10.89 & 10.50 & 0.00 & 0.0399 & 0.1394 \\
$10^{-2}$ & 2.20  & 0.81  & 0.95 & 0.0382 & 0.0884 \\
$0.1$     & 1.10  & 0.12  & 0.00 & 0.0289 & 0.0570 \\
$1.0$     & 0.32  & 0.02  & 0.60 & 0.0314 & 0.0651 \\
\bottomrule
\end{tabular}}
\end{table}

\paragraph{Which source must be resampled?}
Table~\ref{tab:ablation} isolates the two resampling components under protocol (A).  The
baseline resamples only training trajectories, an ``evaluation-only'' interval resamples the
held-out predictions of the trained model, and the joint interval resamples both.  On the
synthetic generator, where the two components are of comparable size, neither single-source
interval is valid ($0.56$ and $0.98$ at $m=100$, $0.85$ and $0.93$ at $m=1000$) and only the
joint interval reaches nominal coverage.  On the real audits the evaluation component dominates,
so the training-only interval collapses to $0.30$ (Adult) and $0.25$ (ACSIncome) --- and to
$0.62$ on COMPAS, where the training component carries a larger share --- while the
evaluation-only interval is close to nominal ($0.97$, $1.00$ and $0.97$).  The component that is
resampled determines coverage and the component that is ignored determines failure.  The
evaluation-only interval is already almost as wide as the joint one ($0.081$ versus $0.082$ on
ACSIncome, $0.098$ versus $0.102$ on Adult, $0.154$ versus $0.173$ on COMPAS), so valid inference
costs little beyond resampling the evaluation set; but only resampling both is uniformly valid
across the two regimes.

\begin{table}[t]
\centering
\small
\setlength{\tabcolsep}{3pt}
\caption{Source ablation under protocol (A), nominal coverage $0.95$: intervals that resample only
training trajectories, only the evaluation set, or both.  Synthetic DI: $T=2\times10^4$, $B=100$,
100 repetitions; real data: $T=3\times10^4$, $B=100$, 60 repetitions.}
\label{tab:ablation}
\resizebox{\linewidth}{!}{\begin{tabular}{lccccccc}
\toprule
setting & $m$ & cov.\ base & cov.\ eval & cov.\ joint & $w$ base & $w$ eval & $w$ joint \\
\midrule
Synthetic (DI) & 100  & 0.56 & 0.98 & 1.00 & 0.107 & 0.270 & 0.283 \\
Synthetic (DI) & 1000 & 0.85 & 0.93 & 0.98 & 0.077 & 0.094 & 0.116 \\
Adult (DI)     & 1000 & 0.30 & 0.97 & 0.98 & 0.025 & 0.098 & 0.102 \\
ACSIncome (DI) & 1000 & 0.25 & 1.00 & 1.00 & 0.015 & 0.081 & 0.082 \\
COMPAS (DM)    & 400  & 0.617 & 0.967 & 1.000 & 0.096 & 0.154 & 0.173 \\
\bottomrule
\end{tabular}}
\end{table}

\section{Discussion}
\label{sec:discussion}

Our results do not contradict the training-side theory of \citet{roy2026uq}; they delimit its scope.
Three practical recommendations follow.  First, auditors should report which estimand an interval
targets, and should not label an interval for the held-out gap as an interval for the deployed
classifier's fairness.  Second, evaluation data should be budgeted before training compute: in our
experiments $m=1000$ labelled observations cannot certify a gap to better than a few percentage
points regardless of $T$.  Third, because the constraint set is itself a random audit design choice,
a single constrained run should be accompanied by a constraint-set sensitivity analysis.

\paragraph{Limitations.}
The theory inherits the asymptotic regime, strong convexity and linear/NTK classifier assumptions of
the baseline.  The joint bootstrap requires $\tbar_T$ to be close to $\theta^*$; unlike the evaluation
and constraint-set components, the optimisation bias is \emph{not} estimable from the same
single-stream bootstrap, and we only report its diagnostic (distance to a directly computed optimum)
and its effect (Table~\ref{tab:sensitivity}).  The constraint-set component assumes that resampling
$\Dp$ is meaningful, that is, that the audit tolerances and group definitions are exchangeable across
audits; when $\Dp$ is a fixed policy object rather than a random audit sample, this component should
be reported as a sensitivity analysis rather than folded into the interval.  Finally, the synthetic
generator supplied by the baseline is close to degenerate (the protected attribute is almost a copy
of the label), so synthetic results should be read as a mechanism check rather than as evidence about
realistic fairness--accuracy trade-offs.  Extensions to multi-class criteria, intersectional group
collections and neural classifiers are left for future work.

\section{Conclusion}
\label{sec:conclusion}
Reported fairness gaps of online-trained classifiers carry three kinds of noise.  We showed that the
existing online bootstrap interval quantifies only the training component, derived the exact rate at
which its coverage degrades as a function of the ratio between training and evaluation effort, and
proposed a joint bootstrap that restores nominal coverage by resampling the training trajectories,
the evaluation set and the constraint set.  A variance budget on three real audits shows that
training noise accounts for less than a tenth of the total uncertainty in typical operating regimes:
evaluation sampling dominates when the held-out set is small, and constraint-set sampling can be the
largest component of all.  The practical message is that fairness audits should budget evaluation
data and audit-design sensitivity before training compute, and should report which estimand their
interval targets.  Future work includes neural classifiers, intersectional group collections and
sequential audits.

\section*{Statements}

\paragraph{Data and code availability.} All datasets are public benchmarks (Adult, ACSIncome and
COMPAS); the processing scripts, data splits, training and bootstrap code, and raw per-repetition
result files are released in the accompanying repository [repository URL / Zenodo DOI to be
inserted].  The reproduction report is included as Appendix~\ref{app:repro}.

\paragraph{Funding and competing interests.} This research received no specific grant from any
funding agency.  The authors declare no competing interests.

\paragraph{Ethics.} The study uses publicly available secondary datasets (Adult, ACSIncome,
COMPAS); no new human-subjects data were collected.

\appendix
\section{Proofs}
\label{app:proofs}

Throughout, $z_j=\mathbb{1}[z=j]$, $P_j=\Pr(z=j)$, $p_j=\Pr(h=1\mid z=j)$ and
$\hat P_j=m^{-1}\sum_{i\in\Dt}z_{i,j}$, $\bar\varphi_j=(\sum_{i\in\Dt}z_{i,j}\varphi_i)/\sum_{i\in\Dt}z_{i,j}$.

\subsection{Proof of Lemma~\ref{lem:eval}}
Write $A=\sum_{i}z_{i,1}\varphi_i$, $B=\sum_i z_{i,1}$, $C=\sum_i z_{i,0}\varphi_i$, $D=\sum_i z_{i,0}$.
Then $\dhat_m(\theta)=A/B-C/D$, $A/m\to P_1p_1$ and $B/m\to P_1$.  Direct computation gives
$\operatorname{Var}(A/m)=P_1p_1(1-P_1p_1)$, $\operatorname{Var}(B/m)=P_1(1-P_1)$ and
$\operatorname{Cov}(A/m,B/m)=P_1p_1(1-P_1)$, and analogously for $(C,D)$.  The delta method with
$J=(1/P_1,-p_1/P_1^2,-1/P_0,p_0/P_0^2)$ yields
\[
v^2(\theta)=\frac{p_1(1-p_1)}{P_1}+\frac{p_0(1-p_0)}{P_0},
\]
which also follows from the law of total variance by conditioning on the realised group memberships.
Independence of $\Dt$ from the training stream and $\Dp$ is part of Assumption~\ref{ass:base}.
For uniformity, the class of half-space indicators is a bounded VC class; on a fixed neighbourhood
of $\theta^*$ it is Donsker, so the empirical process is $O_p(m^{-1/2})$ uniformly in $\theta$.  On
shrinking neighbourhoods, $P(\varphi_\theta\neq\varphi_{\theta^*})\le L\|\theta-\theta^*\|$ for some
$L<\infty$ by the continuity of the feature densities at the boundary, and the same symmetrisation
argument (van der Vaart and Wellner, 1996, Ch.~2.5) gives a local oscillation of order
$o_p(m^{-1/2})+\|\theta-\theta^*\|$ uniformly over $\|\theta-\theta^*\|\le\delta_m$ for any
$\delta_m\to0$; this is what the linearisation step below uses.  The DM case only replaces the
indicator $h=1$ by $h=-1$ and yields a finite sum over the two label strata.

\subsection{Proof of Theorem~\ref{thm:limit}}
Add and subtract the evaluation estimand and split the error into three terms,
\begin{align*}
\sqrt T\big(\dhat_m(\tbar_T)-\Delta(\theta^*)\big)
&=\underbrace{\sqrt T\big(\Delta(\tbar_T)-\Delta(\theta^*)\big)}_{(i)}
+\underbrace{\sqrt T\big(\dhat_m(\theta^*)-\Delta(\theta^*)\big)}_{(ii)}\\
&\quad+\underbrace{\sqrt T\big[(\dhat_m-\Delta)(\tbar_T)-(\dhat_m-\Delta)(\theta^*)\big]}_{(iii)}.
\end{align*}
Term (i) converges to $\mathcal N(0,\sigma^2_\Delta)$ by Lemma~\ref{lem:train}.  For (iii), a
first-order Taylor expansion of the empirical process at $\theta^*$ gives
$(iii)=\sqrt T\,\nabla(\dhat_m-\Delta)(\theta^*)^\top(\tbar_T-\theta^*)+o_p(1)$; the two factors are
$O_p(m^{-1/2})$ and $O_p(T^{-1/2})$ respectively (Lemma~\ref{lem:train}), so
$(iii)=O_p(m^{-1/2})=o_p(1)$ because $m\to\infty$ in the regime $T/m\to\lambda$.  Term (ii) equals
$\sqrt{T/m}\,\sqrt m(\dhat_m(\theta^*)-\Delta(\theta^*))$, which by Lemma~\ref{lem:eval} converges to
$\mathcal N(0,\lambda v^2)$ and is independent of (i) by assumption.  Slutsky's lemma finishes the
proof.

\subsection{Proof of Corollary~\ref{cor:coverage}}
Let $\sigma_{\mathrm{tot}}^2=\sigma^2_\Delta+\lambda v^2$.  The interval of \citet{roy2026uq} has
half-width $z_{\alpha/2}\sigma_\Delta/\sqrt T$, hence
\[
\Pr\{\dhat_m(\tbar_T)\in I_T\}
=\Pr\Big\{\Big|\frac{\sqrt T(\dhat_m(\tbar_T)-\Delta(\theta^*))}{\sigma_{\mathrm{tot}}}\Big|
\le z_{\alpha/2}\frac{\sigma_\Delta}{\sigma_{\mathrm{tot}}}\Big\}+o(1)
=2\Phi\Big(z_{\alpha/2}\frac{\sigma_\Delta}{\sigma_{\mathrm{tot}}}\Big)-1 .
\]
The ratio $\sigma_\Delta/\sigma_{\mathrm{tot}}$ equals one if and only if $\lambda v^2=0$, so the
coverage equals $1-\alpha$ if and only if $m\gg Tv^2/\sigma^2_\Delta$.

\subsection{Proof of Remark~\ref{rem:protocol}}
With a fixed evaluation set, the interval error is
$\dhat_m(\tbar_T)-\Delta(\theta^*)=(\Delta(\tbar_T)-\Delta(\theta^*))+b$ with
$b=\dhat_m(\theta^*)-\Delta(\theta^*)$ fixed conditional on the evaluation set.  By
Lemma~\ref{lem:train},
\[
\Pr\Big\{\big|\sqrt T b+\sqrt T(\Delta(\tbar_T)-\Delta(\theta^*))\big|\le z_{\alpha/2}\sigma_\Delta\Big\}
=\Phi\big(z_{\alpha/2}-\sqrt T b/\sigma_\Delta\big)-\Phi\big(-z_{\alpha/2}-\sqrt T b/\sigma_\Delta\big)+o(1),
\]
which tends to zero for $|b|>0$; the simulations average this expression over the law of $b$, whose
standard deviation is $v/\sqrt m$.  This is the sense in which coverage statements about population
fairness require a fresh evaluation set.

\subsection{Proof of Theorem~\ref{thm:joint}}
Condition on the training data.  For a training multiplier trajectory,
$\sqrt T(\dhat_m(\bar\theta_T^{\,b})-\dhat_m(\tbar_T))=\sqrt T(\Delta(\bar\theta_T^{\,b})-\Delta(\tbar_T))+\sqrt T\big[(\dhat_m-\Delta)(\bar\theta_T^{\,b})-(\dhat_m-\Delta)(\tbar_T)\big]$;
by the online multiplier bootstrap theorem of \citet{roy2026uq} the first term has conditional law
converging to $\mathcal N(0,\sigma^2_\Delta)$, and the second is $o_p(1)$ by the same localisation
argument as in the proof of Theorem~\ref{thm:limit}, because both iterates lie within
$O_p(T^{-1/2})$ of each other.  The evaluation term
$\sqrt m(\dhat_{m,r}(\tbar_T)-\dhat_m(\tbar_T))$ is a standard empirical bootstrap of an i.i.d.\
sample: conditional on the data, the multinomial reweighting gives a conditional law converging to
$\mathcal N(0,v^2)$ by the same delta method as Lemma~\ref{lem:eval}.  The multiplier and evaluation
resamplings use independent randomness, so the two conditional laws are asymptotically independent;
after the rescaling of the main text this gives the law of
$\sqrt T(\dhat_m(\tbar_T)-\Delta(\theta^*))$.  For the third component, each resampled constraint set
yields $\theta^*_s$; when the resampled sets are i.i.d.\ from the audit design, the empirical
distribution over $s$ of $\dhat_m(\theta^*_s)-\dhat_m(\theta^*)$ converges to its sampling law as
$S\to\infty$, provided the constrained optima are estimated with $o_p(S^{-1/2})$ accuracy.  The
(modified) percentile interval is therefore asymptotically valid.  The folded interval for
$|\Delta(\theta^*)|$ has coverage at least that of the signed interval: if the signed interval $[l,u]$
contains $\Delta>0$ then $u\ge\Delta$ and hence $\max(|l|,|u|)\ge|\Delta|$; the case $\Delta<0$ is
symmetric, and when the signed interval contains both signs the folded interval is $[0,\max(|l|,|u|)]$
and covers $|\Delta|$ trivially.

\subsection{Derivation of the sample-size allocation}
Minimising $c_T T+c_m m$ subject to $z^2_{\alpha/2}(\sigma^2_\Delta/T+v^2/m)\le w^2$ has Lagrangian
$c_T T+c_m m+\mu(\sigma^2_\Delta/T+v^2/m-w^2/z^2_{\alpha/2})$.  The first-order conditions give
$T=\sigma_\Delta\sqrt{\mu/c_T}$ and $m=v\sqrt{\mu/c_m}$, hence
$m/T=(v/\sigma_\Delta)\sqrt{c_T/c_m}$.  Substituting into the constraint yields
$\sqrt\mu=z^2_{\alpha/2}(\sigma_\Delta\sqrt{c_T}+v\sqrt{c_m})/w^2$ and therefore
\[
T=\frac{z^2_{\alpha/2}}{w^2}\Big(\sigma^2_\Delta+\sigma_\Delta v\sqrt{c_T/c_m}\Big),
\qquad
m=\frac{z^2_{\alpha/2}}{w^2}\Big(v^2+\sigma_\Delta v\sqrt{c_T/c_m}\Big).
\]

\section{Reproduction details}
\label{app:repro}

The reproduction follows the published equations, with the deviations listed in
Table~\ref{tab:deviations}; the complete report, code and raw result files are released with this
paper.  No official implementation accompanies \citet{roy2026uq}, so the training procedures
(weighted dual averaging with slab projection for DI, penalised SGD for DM) and the online multiplier
bootstrap were re-implemented from the paper.  The most consequential deviations are the step sizes:
in the synthetic DI experiment the published step size $\eta_0=0.01$ leaves the averaged iterate far
from the constrained optimum even at $T=10^5$, so we used $\eta_0=0.1$; the synthetic DM generator
scales the loss by $R^2=500$ and required $\eta_0=0.03$ to avoid divergence.  On the real datasets we
regularised by $\kappa=0.1$ so that the constrained optimum is reached within the training horizon;
this choice is analysed in Table~\ref{tab:sensitivity}.  The complementary sensitivity analysis for
the DM penalty effective at the averaged iterate is reported in the released report.

\begin{table}[t]
\centering
\small
\setlength{\tabcolsep}{4pt}
\caption{Deviations of the reproduction from the published settings (all other settings as in the
original Table~1).  ``z-score'' refers to normalising real features before training.}
\label{tab:deviations}
\resizebox{\linewidth}{!}{\begin{tabular}{lccc}
\toprule
setting & baseline & ours & reason \\
\midrule
step size $\eta_0$ (synthetic DI) & 0.01 & 0.1 & convergence within $T=2\times10^4$ \\
step size $\eta_0$ (synthetic DM) & 0.01 & 0.03 & divergence under $R^2=500$ scaling \\
regularisation $\kappa$ (real) & $10^{-4}$ & 0.1 & reach constrained optimum \\
features (real) & raw & z-score & conditioning \\
repetitions (real) & 100/200 & 100--150 & compute budget \\
$\theta^*$ & trained & direct solve & audit-protocol baseline \\
population $\phi(\theta^*)$ & held-out & full pool & estimand definition \\
\bottomrule
\end{tabular}}
\end{table}

\bibliography{refs}
\end{document}
